\section{Лекция 13}
\subsection{Абсолютная сходимость несобственных интегралов}
Для доказательства дальнейшей теоремы нам понадобится ввести следующие обозначения: Пусть $f$ определена на $\Omega$, обозначим
\[
\begin{aligned}
f^{+}(x) &= \begin{cases}
f(x), & f(x) \geq 0,\\
0,    & f(x) < 0.
\end{cases} &
f^{-}(x) &= \begin{cases}
-f(x), & f(x) \leq 0,\\
0,     & f(x) > 0.
\end{cases}
\end{aligned}
\]
Итак
\[f(x)=f^+(x)-f^-(x),\ \ |f(x)|=f^+(x)+f^-(x)\]
\[f_+(x)=\frac{f(x)+|f(x)|}{2},\ \ f_-(x)=\frac{|f(x)|-f(x)}{2}\]
% тут еще что то про компазицию
Если $E$ --- измеримо по Жордану, то $f\in \mathcal{R}(E)\Rightarrow f^+,f^-, |f|\in \mathcal{R}(E)$.
\begin{exercise}
    Приведите примеры таких $f$ и $E$, что 
    \begin{enumerate}
        \item $f^+\in \mathcal{R}(E),\ f\not\in \mathcal{R}(E)$
        \item $f^-\in \mathcal{R}(E),\ f\not\in \mathcal{R}(E)$
        \item $|f|\in \mathcal{R}(E),\ f\not\in \mathcal{R}(E)$
    \end{enumerate}
\end{exercise}
\noindent Но если хотя бы две из этих функций интегрируемы, то $f$ --- тоже интегрируема.
%смешная картинка
\begin{lemma}
    Пусть $K\subset \Omega$ --- измеримый компакт. Тогда $\forall \epsilon>0\ \exists\ F_{\epsilon}\subset K$ --- измеримый компакт, что 
    \[\idotsint\limits_{F_{\epsilon}}f\ dx_1 \dots dx_k>\idotsint\limits_{K}f^+\ dx_1 \dots dx_k - \epsilon\]
    прчием на $F_{\epsilon}$ выполнено, что $f\equiv f^+$.
\end{lemma}
\begin{proof}
    % \[\idotsint\limits_{K} f^+\ dx_1 \dots dx_k\]
    \[\overline{\overline{s}}(f^+, T)=\sum\limits_{j} \inf\limits_{\Delta_j}f^+\cdot \mu(\Delta_j)\]
    Заметим, что если $\exists\ x\in \Delta_j$, что $f(x)\leq 0$, то $f^+(x)=0$, а значит $\inf\limits_{\Delta_j}f^+=0$. Поэтому достаточно суммировать только по тем $\Delta_j$, на которых принимаюстя положительные значения. Обозначим такие брусы $\Delta_j'$. Тогда
    \[\overline{\overline{s}}(f^+, T)=\sum\limits_{j}\inf\limits_{\Delta_j'} f^+\cdot \mu(\Delta_j')=\sum\limits_{j}\inf\limits_{\Delta_j'} f\cdot \mu(\Delta_j')\]
    Нижняя сумма является инфинумом, поэтому $\forall \epsilon>0\ \exists\ T$:
    \[\overline{\overline{s}}(f^+, T)>\idotsint\limits_{K}f^+\ dx_1 \dots dx_k -\epsilon\]
    Обозначим
    \[F_{\epsilon}=\bigcup\limits_{j} \Delta_j'\] %по положительным
    Пусть $T'$ --- разбиение $F_{\epsilon}$, состоящее из брусов $\Delta_j'$. Тогда
    \[\sum\limits_{j}\inf\limits_{\Delta_j'} f\cdot \mu(\Delta_j')=\overline{\overline{s}}(f,T')\leq \idotsint\limits_{F_{\epsilon}}f\ dx_1\dots dx_k\]
\end{proof}
\begin{theorem}[Критерий интегрируемости в несобственном смысле] $\\$
    Пусть $f\in \mathcal{R}_{\text{loc}}(\Omega)$. Тогда $f\in \widetilde{\mathcal{R}}(\Omega) \Leftrightarrow |f|\in \widetilde{\mathcal{R}}(\Omega)$.
\end{theorem}
\begin{proof}\tab % По Солодову
    \begin{itemize}
        \item[$(\Leftarrow)$:] Пусть $|f|\in \widetilde{\mathcal{R}}(\Omega),\ f^+, f^-\in \mathcal{R}_{\text{loc}}(\Omega)$. Заметим, что на $\Omega$ выполнены
        \[|f|\geq f^+\geq 0,\ \ |f|\geq f^-\geq 0\]
        Тогда по первому признаку сравнения 
        $f^+,f^-\in \widetilde{\mathcal{R}}(\Omega)$. Тогда 
        \[f=f^+-f^-\in \widetilde{\mathcal{R}}(\Omega)\]
        \item[$(\Rightarrow)$:] Предположим, что $f\in \widetilde{\mathcal{R}}(\Omega)$, но $|f|\in \widetilde{\mathcal{R}}(\Omega)$. Это в точности означает, что
        \[\idotsint\limits_{\Omega}|f|\ dx_1 \dots dx_k=+\infty\]
        Отсюда
        \[f^+=\frac{|f|+f}{2}\ \Rightarrow\ \idotsint\limits_{\Omega}f^+\ dx_1 \dots dx_k=+\infty\]
        \[f^-=\frac{|f|-f}{2}\ \Rightarrow\ \idotsint\limits_{\Omega}f^-\ dx_1 \dots dx_k=+\infty\]
        Пусть $\{K_n\}$ --- исчерпание $\Omega$. Хотим построить такое исчерпание $\{K_m'\}$, что при нечетном $m$:
        \[\idotsint\limits_{K_m'}f\ dx_1 \dots dx_k>m\]
        а при четном $m$:
        \[\idotsint\limits_{K_m'}f\ dx_1 \dots dx_k<-m\]
        Построим такое исчерпание по индукции.\\
        \underline{Шаг 1:} Рассмотрим
        \[\idotsint\limits_{K_1}|f|\ dx_1 \dots dx_k\]
        Выберем такое $n_1$, что
        \[\idotsint\limits_{K_{n_1}}f^+\ dx_1 \dots dx_k> \idotsint\limits_{K_1}|f|\ dx_1\dots dx_k+2\]
        По лемме, существует измеримый компакт $Q_1$ такой, что $Q_1\subset K_{n_1}$ и
        \begin{multline*}
            \idotsint\limits_{Q_1}f\ dx_1 \dots dx_k>\idotsint\limits_{K_{n_1}}f^+\ dx_1 \dots dx_k-1>\\
            >\idotsint\limits_{K_1}|f|\ dx_1 \dots dx_k+1
        \end{multline*}
        Определим $K_1'=K_1\cup Q_1$. Тогда
        \begin{multline*}
            \idotsint\limits_{K_1'}f\ dx_1 \dots dx_k=\\
            =\idotsint\limits_{Q_1}f\ dx_1 \dots dx_k+\idotsint\limits_{K_1\setminus Q_1}f\ dx_1 \dots dx_k>\tab[2.5cm]\\
            >\idotsint\limits_{K_1}|f|\ dx_1 \dots dx_k +1 -\idotsint\limits_{K_1}|f|\ dx_1 \dots dx_k=1
        \end{multline*} 
        Второе слагаемое было оценено так:
        \[\idotsint\limits_{K_1\setminus Q_1}f\ dx_1 \dots dx_k\geq -\idotsint\limits_{K_1\setminus Q_1}|f|\ dx_1 \dots dx_k\geq -\idotsint\limits_{K_1}|f|\ dx_1 \dots dx_k\]
        \underline{Шаг 2:} Рассмотрим
        \[\idotsint\limits_{K_2}|f|\ dx_1 \dots dx_k\]
        Выберем такое $n_2>n_1$, что $K_1'\subset K^o_{n_2}$ и
        \[\idotsint\limits_{K_{n_2}}f^-\ dx_1 \dots dx_k > \idotsint\limits_{K_1'\cup K_2}|f|\ dx_1\dots dx_k+3\]
        Применим лемму к функции $-f$. Тогда существует измеримый компакт $Q_2\subset K_{K_{n_2}}$ такой, что
        \begin{multline*}
            \idotsint\limits_{Q_2}(-f)\ dx_1 \dots dx_k > \idotsint\limits_{K_{n_2}}f^-\ dx_1 \dots dx_k -1>\\
            >\idotsint\limits_{K_1'\cup K_2}|f|\ dx_1 \dots dx_k+2
        \end{multline*}
        Определим компакт $F=K_1'\cup Q_2\cup K_2$. Тогда
        \begin{multline*}
            \idotsint\limits_{F}(-f)\ dx_1 \dots dx_k=\\
            =\idotsint\limits_{Q_2}(-f)\ dx_1 \dots dx_k+\idotsint\limits_{(K_1'\cup K_2)\setminus Q_2}(-f)\ dx_1 \dots dx_k>\\
            >\idotsint\limits_{K_1'\cup K_2}|f|\ dx_1 \dots dx_k+2-\idotsint\limits_{K_1'\cup K_2}|f|\ dx_1 \dots dx_k=2
        \end{multline*}
        Второе слагаемое было оценено так:
        \begin{multline*}
            \idotsint\limits_{(K_1'\cup K_2)\setminus Q_2}(-f)\ dx_1 \dots dx_k\geq -\idotsint\limits_{(K_1'\cup K_2)\setminus Q_2}|f|\ dx_1 \dots dx_k\geq\\
            \geq -\idotsint\limits_{K_1'\cup K_2}|f|\ dx_1 \dots dx_k
        \end{multline*}
        По свойству абсолютной непрерывности существует такое $K_2'$, что\\
        $F\subset (K_2')^o$ и тогда
        \[\idotsint\limits_{K_2'} (-f)\ dx_1 \dots dx_k>2\]
        % Хотим построить $K'_{2j+1}$. По предположению индукции построено $K_{2j}'$
        % \[\idotsint\limits_{K_{2j}'}f\ dx_1 \dots dx_k<-2j\]
        Шаг индукции. Пусть построены все $K_m'$ при $m<2j$. \\
        \underline{Шаг $2j$:} Рассмотрим
        \[\idotsint\limits_{K_{2j}}|f|\ dx_1 \dots dx_k\]
        Выберем такое $n_{2j}>n_{2j-1}$, что $K_{2j-1}'\subset K^o_{n_{2j}}$ и
        \[\idotsint\limits_{K_{n_{2j}}}f^-\ dx_1 \dots dx_k > \idotsint\limits_{K_{2j-1}'\cup K_{2j}}|f|\ dx_1\dots dx_k+(2j+1)\]
        Применим лемму к функции $-f$. Тогда существует измеримый компакт $Q_{2j}\subset K_{n_{2j}}$ такой, что
        \begin{multline*}
            \idotsint\limits_{Q_{2j}}(-f)\ dx_1 \dots dx_k > \idotsint\limits_{K_{n_{2j}}}f^-\ dx_1 \dots dx_k -1 >\\
            >\idotsint\limits_{K_{2j-1}'\cup K_{2j}}|f|\ dx_1 \dots dx_k+2j
        \end{multline*}
        Определим компакт $G=K_{2j-1}'\cup Q_{2j}\cup K_{2j}$. Тогда
        \begin{multline*}
            \idotsint\limits_{G}(-f)\ dx_1 \dots dx_k = \\
            =\idotsint\limits_{Q_{2j}}(-f)\ dx_1 \dots dx_k + \idotsint\limits_{(K_{2j-1}'\cup K_{2j})\setminus Q_{2j}}(-f)\ dx_1 \dots dx_k > \\
            >\idotsint\limits_{K_{2j-1}'\cup K_{2j}}|f|\ dx_1 \dots dx_k+2j-\idotsint\limits_{K_{2j-1}'\cup K_{2j}}|f|\ dx_1 \dots dx_k = 2j
        \end{multline*}
        Второе слагаемое было оценено так:
        \begin{multline*}
            \idotsint\limits_{(K_{2j-1}'\cup K_{2j})\setminus Q_{2j}}(-f)\ dx_1 \dots dx_k \geq -\idotsint\limits_{(K_{2j-1}'\cup K_{2j})\setminus Q_{2j}}|f|\ dx_1 \dots dx_k \geq \\
            \geq-\idotsint\limits_{K_{2j-1}'\cup K_{2j}}|f|\ dx_1 \dots dx_k
        \end{multline*}
        По свойству абсолютной непрерывности существует такое $K_{2j}'$, что \\
        $G\subset (K_{2j}')^o$ и тогда
        \[\idotsint\limits_{K_{2j}'}(-f)\ dx_1 \dots dx_k > 2j\]
        \underline{Шаг $2j+1$:}
        Рассмотрим
        \[\idotsint\limits_{K_{2j+1}}|f|\ dx_1 \dots dx_k\]
        Выберем такое $n_{2j+1}>n_{2j}$, что $K_{2j}'\subset K^o_{n_{2j+1}}$ и
        \[\idotsint\limits_{K_{n_{2j+1}}}f^+\ dx_1 \dots dx_k > \idotsint\limits_{K_{2j}'\cup K_{2j+1}}|f|\ dx_1\dots dx_k+(2j+2)\]
        По лемме, существует измеримый компакт $Q_{2j+1}\subset K_{n_{2j+1}}$ такой, что
        \begin{multline*}
            \idotsint\limits_{Q_{2j+1}}f\ dx_1 \dots dx_k > \idotsint\limits_{K_{n_{2j+1}}}f^+\ dx_1 \dots dx_k -1 > \\ 
        \idotsint\limits_{K_{2j}'\cup K_{2j+1}}|f|\ dx_1 \dots dx_k+(2j+1)
        \end{multline*}
        Определим компакт $S=K_{2j}'\cup Q_{2j+1}\cup K_{2j+1}$. Тогда
        \begin{multline*}
        \idotsint\limits_{S}f\ dx_1 \dots dx_k=\\
        =\idotsint\limits_{Q_{2j+1}}f\ dx_1 \dots dx_k+\idotsint\limits_{(K_{2j}'\cup K_{2j+1})\setminus Q_{2j+1}}f\ dx_1 \dots dx_k > \\
        >\idotsint\limits_{K_{2j}'\cup K_{2j+1}}|f|\ dx_1 \dots dx_k+(2j+1) -
        \idotsint\limits_{K_{2j}'\cup K_{2j+1}}|f|\ dx_1 \dots dx_k = 2j+1
        \end{multline*}
        Второе слагаемое было оценено так:
        \begin{multline*}
            \idotsint\limits_{(K_{2j}'\cup K_{2j+1})\setminus Q_{2j+1}}f\ dx_1 \dots dx_k \geq-\idotsint\limits_{(K_{2j}'\cup K_{2j+1})\setminus Q_{2j+1}}|f|\ dx_1 \dots dx_k \geq\\
            \geq-\idotsint\limits_{K_{2j}'\cup K_{2j+1}}|f|\ dx_1 \dots dx_k
        \end{multline*}
        По свойству абсолютной непрерывности существует такое $K_{2j+1}'$, что \\
        $S\subset (K_{2j+1}')^o$ и тогда
        \[\idotsint\limits_{K_{2j+1}'}f\ dx_1 \dots dx_k > 2j+1\]
    \end{itemize}
\end{proof}
\begin{example}
    Сферические координаты в многомерном пространстве.
\end{example}